% AST2000 Part 5: planning and flying the journey to the destination planet.
% Compile: cd "Project 5/paper" && latexmk -pdf part5.tex   (or: python build_papers.py 5 from the repo root)
% Every number written as \val{key} comes from results.tex, made by pro_5.py.
\documentclass[aps,prd,twocolumn,nofootinbib,floatfix,superscriptaddress]{revtex4-2}
\usepackage{../../common/paper/ast2000}

\begin{document}

\title{Part 5: Planning and Flying the Journey to Planet \val{dest}}
\author{Sakarias Frette}
\affiliation{AST2000 Rerun, University of Oslo}
\date{\today}

\begin{abstract}
We plan and fly a journey from the home planet into a circular orbit around planet \val{dest}.
A trajectory simulator with exact planetary ephemerides and an adaptive eighth-order integrator reproduces a 20-orbit test to \val{kepler-err-km}\,km.
A porkchop plot of $\val{grid}^2$ Lambert transfers selects a launch at $t=\val{t-launch}$\,yr and a \val{travel-days}-day trip.
The plan has two key elements:
\begin{itemize}
\item a departure burn of \val{plan-boost}\,m/s right after launch, deep in the home planet's gravity well, which replaces \val{vinf-dep}\,km/s far away;
\item an orbit radius chosen to minimise the capture burn.
\end{itemize}
Newton shooting in the full $N$-body simulation refines the plan.
We then flew it with \code{ast2000tools}, using the Part 4 software to find the spacecraft's state from its instruments.
\val{n-corrections-word} course corrections costing \val{correction-total}\,cm/s in total brought the spacecraft to \val{closest}\,km from planet \val{dest}, where a \val{real-injection}\,m/s burn put it into a circular orbit with a period of \val{orbit-P}\,h.
The whole trip used \val{real-total}\,km/s and left \val{fuel-final}\,kg of fuel for the landing.
\end{abstract}

\maketitle

% =====================================================================
\section{Introduction}
Part 3 chose planet \val{dest} as the destination and estimated the cost of the trip.
This part~\cite{ast2000project} makes the trip, in four steps:
\begin{itemize}
\item a trajectory simulator (\cref{sec:sim});
\item a plan that minimises fuel (\cref{sec:plan});
\item the real flight, which can only learn its state from its instruments (\cref{sec:flight});
\item a description of the orbit we end in (\cref{sec:orbit}).
\end{itemize}
Solar system seed: \val{seed}.

% =====================================================================
\section{The trajectory simulator}\label{sec:sim}
While coasting, the spacecraft feels the star and all the planets,
\begin{equation}
  \ddot{\vb r} = -\frac{GM_\star\vb r}{|\vb r|^3}-\sum_i\frac{GM_i(\vb r-\vb r_i(t))}{|\vb r-\vb r_i(t)|^3}.
\end{equation}

\paragraph{Ephemeris.}
The planets' positions come from Kepler's equation (Part 2) on a uniform grid every \val{eph-dt-min} minutes.
Between grid points we interpolate with cubic Hermite polynomials, which also match the exact velocities.
\begin{itemize}
\item The interpolation error is \val{eph-err-hermite}\,m, against \val{eph-err-linear}\,m for linear interpolation.
\item The planets' velocities stay continuous. With linear interpolation they jump at every grid point, and an adaptive integrator wastes steps on these artificial kicks: \val{linear-evals} force evaluations against \val{hermite-evals} for the same trip.
\end{itemize}

\paragraph{Integrator.}
The trip has two very different time scales: minutes near a planet and months in between.
We therefore use the adaptive DOP853 integrator~\cite{dormand1980} with a relative tolerance of \val{rtol}.
A test particle on a Kepler orbit stays within \val{kepler-err-km}\,km of the exact solution after 20 orbits (\cref{fig:sim}).

\paragraph{Batches.}
The simulator integrates many spacecraft at once, computing the gravity of all planets on all spacecraft with NumPy broadcasting.
A batch of \val{batch-max} costs \val{batch-speedup} times less per spacecraft than single runs.
The planner gets its Jacobians this way.

\widefig{a_simulator_tests}{Simulator tests: ephemeris interpolation error, accuracy and cost against tolerance, and cost per spacecraft in a batch.}{fig:sim}

% =====================================================================
\section{Planning}\label{sec:plan}
\subsection{When to go: the porkchop plot}
For every pair of launch date and travel time on a $\val{grid}\times\val{grid}$ grid, Lambert's problem~\cite{izzo2015} gives the transfer orbit around the star.
Its velocities relative to the planets at the two ends, $v_\infty$, set the cost.

\paragraph{Departure.}
The Part 1 launch ends at exactly the escape speed $v_{\rm esc}$, still deep in the home planet's gravity well.
A burn there gains much more energy than the same burn far away (the Oberth effect).
Leaving with $v_\infty$ costs only
\begin{equation}
  \Delta v_{\rm dep}=\sqrt{v_\infty^2+v_{\rm esc}^2}-v_{\rm esc}.
\end{equation}

\paragraph{Arrival.}
Falling in from $v_\infty$ to a circular orbit of radius $r$ costs
\begin{equation}
  \Delta v_{\rm arr}=\sqrt{v_\infty^2+\frac{2GM}{r}}-\sqrt{\frac{GM}{r}},
\end{equation}
which is smallest at $r=2GM/v_\infty^2$, where $\Delta v_{\rm arr}=v_\infty/\sqrt2$.
So the planner chooses the orbit radius as well as the route.

\paragraph{Result.}
\Cref{fig:porkchop} shows windows once per synodic period.
The cheapest point in $\val{window-lo}\le t\le\val{window-hi}$\,yr costs \val{porkchop-min}\,km/s in total: launch at $t=\val{t-launch}$\,yr and arrive \val{travel-days} days later.
There, $v_\infty=\val{vinf-dep}$\,km/s at departure and \val{vinf-arr}\,km/s at arrival, so the burns are \val{pc-dv-dep} and \val{pc-dv-arr}\,m/s.
\Cref{fig:trades} shows both trades.
The Oberth departure saves most of the departure cost, and the chosen orbit, \val{r-orbit-radii} planet radii, sits at the minimum of the capture curve.

\subsection{Where to launch from, and shooting}
The launch goes straight up from the launch site (Part 3).
We choose the site whose ``up'' points along the required departure direction ($\phi=\val{launch-angle}^\circ$).
The rocket then leaves the planet in a straight line: gravity slows it but cannot bend it (\cref{fig:departure}).

The patched-conic numbers are only a first guess, off by \val{guess-off}\,m/s.
We refine the departure burn with Newton's method in the full $N$-body simulation.
For small changes the miss at arrival is linear in the burn, and one batched run of three spacecraft gives the $2\times2$ Jacobian.
An outer secant loop adjusts the aim point beside the planet until the closest approach equals the planned orbit radius.

The final plan is:
\begin{itemize}
\item a departure burn of \val{plan-boost}\,m/s;
\item closest approach after \val{arrival-days} days at \val{r-orbit}\,km;
\item an injection burn of \val{plan-inj}\,m/s;
\item fuel of \val{plan-fuel-dep}\,kg and \val{plan-fuel-inj}\,kg for the two burns, out of \val{fuel-after-launch}\,kg (\cref{fig:budget}).
\end{itemize}

\subsection{Sensitivity}
\Cref{fig:sens} shows how far the arrival point moves per m/s of velocity error, depending on when the error happens:
\begin{itemize}
\item right after launch: \val{sens-start}\,km;
\item halfway: \val{sens-mid}\,km;
\item a day before arrival: \val{sens-end}\,km.
\end{itemize}
Two conclusions follow.
The departure burn must be exact, and course corrections are cheapest early.

\widefig{b1_porkchop}{Porkchop plot of the total $\Delta v$ against launch date and travel time (left), and the cost of the best travel time for every launch date (right).}{fig:porkchop}
\widefig{b7_trades}{Left: capture burn against orbit radius, with the optimum $2GM/v_\infty^2$. Right: departure burn against the $v_\infty$ needed, with and without the Oberth effect.}{fig:trades}
\widefig{b2_planned_journey}{The planned trip in the star's frame, and the distance to planet \val{dest} against time.}{fig:plan}
\widefig{b3_departure}{The first 48 hours in the home planet's frame: launching ``up'' along the departure direction.}{fig:departure}
\widefig{b4_arrival}{Arrival in planet \val{dest}'s frame: the approach and the circular orbit after injection.}{fig:arrival}
\widefig{b5_budget}{The planned burns and the fuel left after each.}{fig:budget}
\widefig{b6_sensitivity}{Miss at arrival per 1\,m/s velocity error, against the time of the error.}{fig:sens}

% =====================================================================
\section{The flight}\label{sec:flight}
\subsection{Method}
We launched at the planned time and site (verification: \val{launch-verified}) and ran the Part 4 orientation software (verification: \val{orient-verified}).
Our launch prediction differs from the API's by a few m/s, which a four-month trip would amplify.
We therefore re-aimed the departure burn from the \emph{measured} state, giving \val{real-departure}\,m/s.

During the cruise the spacecraft oriented itself at \val{n-correction-points} points.
At each it predicted where it would arrive with our simulator, and corrected only if it was heading more than \val{miss-tolerance}\,km off the aim point.
At the end it predicted the closest approach, coasted there, and burned into a circular orbit.

\subsection{Results}
The first check, a few days out, found us heading \val{first-off}\,km off.
Only \val{n-corrections} corrections were needed, together \val{correction-total}\,cm/s, and the last check found us \val{last-off}\,km off (\cref{fig:corr}).
The closest approach came at $t=\val{t-closest}$\,yr, \val{closest}\,km from the planet's centre, which is exactly the planned radius.
The injection burn was \val{real-injection}\,m/s.

The telemetry (\cref{fig:telemetry}) shows the physics of the trip.
\begin{itemize}
\item The speed relative to the star falls as the spacecraft climbs outward.
\item The energy relative to the star stays constant through the cruise. It changes only while the planets' gravity dominates: climbing out of the home planet's well and falling into the destination's.
\item The fuel steps down at the two large burns; the corrections do not show at this scale.
\end{itemize}
The total cost was \val{real-total}\,km/s, against \val{plan-total}\,km/s planned, leaving \val{fuel-final}\,kg of fuel.

\widefig{c1_flight}{The real flight: trajectory with burns, deviation from the planned path, and all burns.}{fig:flight}
\widefig{c2_telemetry}{Flight telemetry from the spacecraft's own orientations: range to both planets, speed and orbital energy relative to the star, and fuel. Dotted lines mark burns.}{fig:telemetry}
\widefig{c3_corrections}{Course-correction log: predicted miss of the aim point at each check, and the correction burns.}{fig:corr}

% =====================================================================
\section{The orbit}\label{sec:orbit}
From one orientation relative to the planet we get $r$, $v_r$ and $v_\theta$.
The two-body formulas~\cite{ast2000notes} then give the orbit:
\begin{gather}
  h=rv_\theta,\quad E=\tfrac12v^2-\frac{GM}{r},\quad a=-\frac{GM}{2E},\nonumber\\
  e=\sqrt{1+\frac{2Eh^2}{(GM)^2}},\quad P=2\pi\sqrt{\frac{a^3}{GM}}.
\end{gather}
Right after injection, $a=\val{orbit-a}$\,km.
The periapsis and apoapsis are both \val{orbit-peri}\,km, and the period is \val{orbit-P}\,h.
Over \val{n-orbits-check} orbits, $a$ varies by $\val{orbit-a-var}$ (relative), and $e$ stays below \val{orbit-e} (\cref{fig:orbit}).

After a coast of \val{long-coast}\,yr the spacecraft is still in orbit, between \val{long-peri} and \val{long-apo}\,km.
That orbit has shrunk by \val{api-shrink}\%.
Our own simulator keeps $a$ constant over the same time.
The shrinking comes from the API's RK45 integrator, which is not symplectic and slowly loses energy~\cite{hairer2006}.
It is not real physics.

\widefig{d_orbit}{The orbit around planet \val{dest}. Left: orientations over \val{n-orbits-check} orbits and over \val{long-coast}\,yr. Middle: orbital elements from each orientation. Right: long-term evolution in the API and in our simulator.}{fig:orbit}

% =====================================================================
\section{Discussion and conclusion}
Three insights carried the trip.
\begin{itemize}
\item \textbf{The Oberth effect.} Burning right after launch made the departure cost \val{plan-boost}\,m/s instead of \val{vinf-dep}\,km/s.
\item \textbf{Choosing the orbit radius.} Picking it at $2GM/v_\infty^2$ minimised the capture burn.
\item \textbf{Measuring, not trusting.} Re-aiming from the measured state and correcting early kept the corrections at the cm/s level.
\end{itemize}
The total of \val{real-total}\,km/s compares with \val{hohmann-dv}\,km/s for Part 3's simple Hohmann estimate.

The approximations are as follows.
\begin{itemize}
\item The burns are instantaneous. The real ones would take seconds to minutes, which is short compared with the trip.
\item The planets are on Kepler orbits. This is also what the API assumes, so our simulation matches its flight.
\end{itemize}
The spacecraft is now in orbit around planet \val{dest} with \val{fuel-final}\,kg of fuel, ready for Part 6.

\bibliography{../../common/paper/references}

\appendix
\section{Mission log}
\begin{missionbox}{Part 5 key numbers (seed \val{seed})}
Launch & $t=\val{t-launch}$\,yr, site $\val{launch-angle}^\circ$ \\
Travel time & \val{arrival-days}\,days \\
Departure / injection & \val{real-departure} / \val{real-injection}\,m/s \\
Corrections & \val{n-corrections}, total \val{correction-total}\,cm/s \\
Orbit radius / period & \val{orbit-a}\,km / \val{orbit-P}\,h \\
Total $\Delta v$ & \val{real-total}\,km/s \\
Fuel left & \val{fuel-final}\,kg \\
Launch / orientation check & \val{launch-verified} / \val{orient-verified} \\
\end{missionbox}

\end{document}
